% =============================================================================
% Appendix N: Implementation Infrastructure Details
% =============================================================================

\section{Implementation Infrastructure Details}\label{app:implementation}

This appendix provides the detailed implementation descriptions that were summarised in Section~\ref{sec:architecture}. Each subsection corresponds to a subsystem that is referenced but not fully described in the main text.

\subsection{Calibration Layer}\label{app:calibration}

The calibration layer (\texttt{calibration.py}) mitigates overconfidence in the default $\gamma(C)$ aggregation by offering four complementary strategies: \textbf{linear calibration} (per-actor accuracy-weighted confidence), \textbf{EWMA time-decay} (exponentially weighted moving average with smoothing factor $\alpha = 0.3$), \textbf{epistemic context} (domain-specific accuracy scores per actor per domain), and \textbf{per-band correction} (low-confidence values $[0.0, 0.3)$ receive $+0.15$ boost; high-confidence $[0.7, 1.0]$ receive $-0.10$ dampening). These strategies operate independently and can be combined. \texttt{VALIDATE} events trigger an EWMA accuracy calibration feedback loop: when a \texttt{VALIDATE} event's \texttt{triggers\_transition} field contains \texttt{"validated"}, the actor's accuracy is updated via \texttt{calibrator.update\_accuracy\_ewma()} using the confidence value from the \texttt{VALIDATE} event.

\subsection{Semantic Web Interoperability}\label{app:semantic-web}

ADL Lite implements two Semantic Web export pathways. \textbf{Bidirectional OWL~2 DL import/export} (\texttt{owl\_export.py} / \texttt{owl\_import.py}) converts \texttt{ADLDocument} to OWL~2 DL in Turtle or RDF/XML, mapping concepts to \texttt{NamedIndividual}s, status to annotation properties, and events to typed individuals. Round-trip validation preserves concept identity, status, confidence, and event count. \textbf{RDF-star / SPARQL-star} (\texttt{rdfstar\_export.py}) converts events to embedded triples, enabling provenance queries such as ``find all concepts validated by agent X with confidence $>0.8$ in the last 30 days.''

\subsection{Split-Lock Concurrency Model}\label{app:split-lock}

As of v0.6.0-alpha, the \texttt{EventChain} class employs a dual-lock concurrency design. Two independent \texttt{threading.RLock} instances protect distinct shared state: (i)~\texttt{\_events\_lock} guards the append-only \texttt{\_events} list, and (ii)~\texttt{\_cache\_lock} guards the CRDT-derived caches (\texttt{\_cached\_status}, \texttt{\_cached\_confidence}) and the incremental verification cursor (\texttt{\_verified\_up\_to}). The strict acquisition order (\texttt{\_events\_lock} $\to$ \texttt{\_cache\_lock}) prevents deadlocks. This separation allows status and confidence queries (which only need \texttt{\_cache\_lock}) to proceed concurrently with event appends (which require both locks in sequence).

The \texttt{ConsensusEngine} (\texttt{consensus.py}) supports dynamic mode switching between development and production environments. In dev mode, the effective minimum validator threshold $N_{\min}$ defaults to 1, enabling single-actor testing workflows. In production mode, $N_{\min}$ is set to $\max(\text{ontology\_min}, 2)$, enforcing at least two validators for status transitions.

\subsection{Vector Index for Semantic Search}\label{app:vector-index}

The \texttt{vector\_index.py} module provides a FAISS-backed approximate nearest-neighbour (ANN) index for semantic search over concept embeddings. Each concept's textual content (name, description, L2 body) is embedded via a configurable backend (\texttt{embeddings.py}, supporting SentenceTransformer or OpenAI APIs) and stored alongside metadata in a SQLite database. The index supports batch insertion, top-$k$ ANN retrieval with configurable cosine similarity threshold, and atomic save/load of FAISS index files and SQLite metadata. The vector index is optional (requires the \texttt{[embeddings]} extra) and is used primarily by the canonicalization pipeline to cluster near-duplicate concepts before proposing merges.

\subsection{Merkle Transparency Anchors}\label{app:merkle}

The \texttt{merkle.py} module implements SHA-256 Merkle trees over event hashes, providing $O(\log n)$ inclusion proofs for batch integrity verification. A \texttt{MerkleTree} is constructed from a list of leaf hashes (typically event \texttt{hash} values from an EventChain), producing a root digest that can be anchored to external transparency logs or blockchain-based timestamping services. The module provides \texttt{MerkleTree} (builds the tree and generates inclusion proofs), \texttt{MerkleTree.verify\_proof} (statically verifies an inclusion proof), and \texttt{compute\_chain\_merkle\_root} (convenience function for computing the root over a list of hex hashes). Merkle anchors complement the linear hash chain: while the chain provides sequential integrity (each event linked to its predecessor), the Merkle tree enables efficient batch verification and external anchoring without requiring full-chain re-verification.

\subsection{LLM-Driven Canonicalization}\label{app:canonicalization}

The \texttt{canonicalization.py} module implements a pipeline for detecting and merging near-duplicate concepts using LLM-assisted normalization. The workflow is: (i)~\textbf{Clustering:} the vector index groups concepts by embedding similarity above a configurable threshold. (ii)~\textbf{LLM Proposal:} for each cluster, an LLM backend (pluggable via the \texttt{LLMBackend} protocol; implementations include OpenAI \texttt{OpenAILLMBackend} and Anthropic \texttt{AnthropicLLMBackend}) proposes a canonical form, including a canonical \texttt{adl\_id}, multilingual name, and a set of \texttt{isomorphic-to} relations. (iii)~\textbf{Action Generation:} the engine translates the LLM proposal into auditable ADL action blocks (\texttt{REGISTER} for the canonical concept, \texttt{RELATE} for isomorphic links, optionally \texttt{DEPRECATE} for duplicates). (iv)~\textbf{Dry-Run Safety:} by default, the pipeline runs in dry-run mode (\texttt{dry\_run=True}), generating actions without mutating any EventChains. Callers must explicitly execute the actions after review.

\subsection{REST API and Authentication}\label{app:rest-api}

The \texttt{api.py} module (481 lines) exposes a FastAPI-based REST API providing programmatic access to ADL Lite operations. The API implements 10 endpoints under the \texttt{/api/v1/consensus/} prefix: \texttt{register}, \texttt{transition}, \texttt{status}, \texttt{history}, \texttt{fork}, \texttt{verify}, \texttt{list}, \texttt{mode/dev}, \texttt{mode/production}, and \texttt{mode query}. Authentication supports dual mechanisms: JWT bearer tokens and API key authentication via \texttt{api\_auth.py} (344 lines), enabling flexible integration with both human operators and automated agents. The API enforces CORS policies, rate limiting, and pagination for list endpoints.

\subsection{MCP Server for Agent Integration}\label{app:mcp-server}

The \texttt{mcp\_server.py} module (458 lines) implements a Model Context Protocol (MCP) server using the FastMCP framework, enabling ADL Lite to serve as a governance backend for MCP-compatible LLM agents. The server exposes 10 tools (\texttt{adl\_parse}, \texttt{adl\_validate}, \texttt{adl\_register}, \texttt{adl\_transition}, \texttt{adl\_status}, \texttt{adl\_verify}, \texttt{adl\_history}, \texttt{adl\_fork}, \texttt{adl\_list}, \texttt{adl\_ontology\_query}), 2 resources (\texttt{adl://ontology}, \texttt{adl://capability/\{adl\_id\}}), and 1 prompt (\texttt{adl\_lifecycle\_prompt}). Two transport modes are supported: stdio for local agent integration and streamable-http for networked deployments.

\subsection{SHACL Validation Framework}\label{app:shacl-framework}

The \texttt{shacl\_validation.py} module (334 lines) implements W3C SHACL shape validation using the pyshacl library, providing machine-checkable conformance checking for ADL documents. Six SHACL shapes are defined: \texttt{EventShape}, \texttt{ConceptShape}, \texttt{AgentShape}, \texttt{CalibrateEventShape}, \texttt{RelationShape}, and \texttt{ForkShape}. These shapes validate exported RDF against the structural and semantic constraints described in Appendix~\ref{app:shacl}. The module provides two primary validation functions: \texttt{validate\_adl\_rdf()} for validating RDF graphs and \texttt{validate\_adl\_document()} for direct document validation.

\subsection{Neo4j Graph Adapter}\label{app:neo4j}

The \texttt{neo4j\_adapter.py} module (148 lines) provides a pluggable graph backend implementing the same interface as the NetworkX DiGraph for the WarmIndex. It supports graph operations including \texttt{add\_edge}, \texttt{bfs}, \texttt{\_\_contains\_\_}, \texttt{node\_count}, and \texttt{close}. The adapter supports dual-write mode, maintaining synchronized SQLite and Neo4j stores for query flexibility, and provides a \texttt{rebuild\_from\_relations} method for graph reconstruction. Configuration is provided via environment variables, enabling deployment-specific tuning without code changes.

\subsection{CRDT Multi-Way Merge Generalization}\label{app:crdt-multiway}

The \texttt{merge\_event\_chains(*chains)} function has been generalized to support N$\geq$3 chain merging via \texttt{functools.reduce} over the \texttt{\_merge\_two\_chains} helper, maintaining full backward compatibility with the 2-chain case. This enables concurrent merge of multiple branches without pairwise reduction, improving both performance and semantic clarity for multi-agent scenarios. New proof functions \texttt{prove\_multi\_way\_merge()} and \texttt{prove\_event\_chain\_multi\_way\_merge()} establish CRDT convergence properties (Theorem~9) for arbitrary numbers of concurrent branches. The multi-way merge preserves the established properties: commutability, associativity, idempotence, and deterministic state derivation.
